% GENERATED FILE. Edit canonical lessons, then run npm run generate:books.
% canonical-source-hash: fnv1a64:aeb598af153e90f8
% canonical-lessons: MW-C69-taro, MW-C69-suraj, MW-C69-chad, MW-C69-samandar, MW-C69-ag

\chapter{Star, Sun, Moon, Sea, Fire}
\label{ch:mw-star-sun-moon-sea-fire}
\hlchaptermodality{69}

\begin{chapteropening}
\textbf{By the end of this chapter:} \emph{I can say a star, the sun, the moon, the sea and fire.}

Complete the last lesson of chapter 69: I can say a star, the sun, the moon, the sea and fire.
\end{chapteropening}

\section[\texorpdfstring{\mw{तारो} (tāro)}{तारो (tāro)}]{\mw{तारो} (tāro) — a star}

\label{lesson:MW-C69-taro}



\begin{quote}
Before the new one: say the Marwadi for sand, then the Marwadi for a field.
\end{quote}

\subsection*{You'll want to know: \mw{तारो}}
\textbf{\mw{तारो}} — \emph{tāro} — \textquotedblleft{}a star\textquotedblright{}.

\begin{cousinweb}
\textbf{\mw{तारो}} (\emph{tāro}) is a tadbhava, a word inherited from Sanskrit and reshaped by sound change. It comes from \emph{tārā}, \textquotedblleft{}fixed star\textquotedblright{}. CDIAL gives no Marwari form in this entry (CDIAL 5798); the match is through Hindi \emph{tārā} and Gujarati \emph{tāro}. English \emph{star} is a distant cousin.
\end{cousinweb}

\begin{grammarlens}[title={What you've built}]
One more word for this chapter.
\end{grammarlens}

\subsection*{Your turn}
\begin{itemize}
\raggedright
  \item \emph{Say it:} \emph{tāro}
  \item \emph{Say it:} \emph{tāro}, once more
  \item \emph{Say it:} say \emph{tāro}
  \item \emph{Recall:} say the Marwadi for sand, then the Marwadi for a field, then say \emph{tāro} again
\end{itemize}

\begin{tcolorbox}[breakable,colback=teal!4,colframe=teal!35!black,title={Before you move on}]
What does \mw{तारो} mean? (\textquotedblleft{}A star\textquotedblright{}.) Say it once more.
\end{tcolorbox}
\section[\texorpdfstring{\mw{सूरज} (sūraj)}{सूरज (sūraj)}]{\mw{सूरज} (sūraj) — the sun}

\label{lesson:MW-C69-suraj}



\begin{quote}
Before the new one: say the Marwadi for a field, then the Marwadi for a star.
\end{quote}

\subsection*{You'll want to know: \mw{सूरज}}
\textbf{\mw{सूरज}} — \emph{sūraj} — \textquotedblleft{}the sun\textquotedblright{}.

\begin{cousinweb}
\textbf{\mw{सूरज}} (\emph{sūraj}) is usually traced to Sanskrit \emph{sūrya}, \textquotedblleft{}sun\textquotedblright{}. CDIAL lists the Sanskrit word (CDIAL 13574) but no Marwari or Hindi form of it, so the exact path is not settled.
\end{cousinweb}

\begin{grammarlens}[title={What you've built}]
One more word for this chapter.
\end{grammarlens}

\subsection*{Your turn}
\begin{itemize}
\raggedright
  \item \emph{Say it:} \emph{sūraj}
  \item \emph{Say it:} \emph{sūraj}, once more
  \item \emph{Say it:} say \emph{sūraj}
  \item \emph{Recall:} say the Marwadi for a field, then the Marwadi for a star, then say \emph{sūraj} again
\end{itemize}

\begin{tcolorbox}[breakable,colback=teal!4,colframe=teal!35!black,title={Before you move on}]
What does \mw{सूरज} mean? (\textquotedblleft{}The sun\textquotedblright{}.) Say it once more.
\end{tcolorbox}
\section[\texorpdfstring{\mw{चांद} (chā̃d)}{चांद (chā̃d)}]{\mw{चांद} (chā̃d) — the moon}

\label{lesson:MW-C69-chad}



\begin{quote}
Before the new one: say the Marwadi for a star, then the Marwadi for the sun.
\end{quote}

\subsection*{You'll want to know: \mw{चांद}}
\textbf{\mw{चांद}} — \emph{chā̃d} — \textquotedblleft{}the moon\textquotedblright{}.

\begin{cousinweb}
\textbf{\mw{चांद}} (\emph{chā̃d}) is a tadbhava, a word inherited from Sanskrit and reshaped by sound change. It comes from \emph{candra}, \textquotedblleft{}shining; moon\textquotedblright{}. CDIAL gives no Marwari form in this entry (CDIAL 4661); the match is through Hindi \emph{cā̃d} and Gujarati \emph{cā̃d}.
\end{cousinweb}

\begin{grammarlens}[title={What you've built}]
One more word for this chapter.
\end{grammarlens}

\subsection*{Your turn}
\begin{itemize}
\raggedright
  \item \emph{Say it:} \emph{chā̃d}
  \item \emph{Say it:} \emph{chā̃d}, once more
  \item \emph{Say it:} say \emph{chā̃d}
  \item \emph{Recall:} say the Marwadi for a star, then the Marwadi for the sun, then say \emph{chā̃d} again
\end{itemize}

\begin{tcolorbox}[breakable,colback=teal!4,colframe=teal!35!black,title={Before you move on}]
What does \mw{चांद} mean? (\textquotedblleft{}The moon\textquotedblright{}.) Say it once more.
\end{tcolorbox}
\section[\texorpdfstring{\mw{समंदर} (samandar)}{समंदर (samandar)}]{\mw{समंदर} (samandar) — the sea}

\label{lesson:MW-C69-samandar}



\begin{quote}
Before the new one: say the Marwadi for the sun, then the Marwadi for the moon.
\end{quote}

\subsection*{You'll want to know: \mw{समंदर}}
\textbf{\mw{समंदर}} — \emph{samandar} — \textquotedblleft{}the sea\textquotedblright{}.

\begin{cousinweb}
\textbf{\mw{समंदर}} (\emph{samandar}) goes back to Sanskrit \emph{samudra}, \textquotedblleft{}sea\textquotedblright{}, with an added nasal, like the Old Awadhi \emph{samuṃda} that CDIAL lists under it (CDIAL 13236). CDIAL gives no Marwari form.
\end{cousinweb}

\begin{grammarlens}[title={What you've built}]
One more word for this chapter.
\end{grammarlens}

\subsection*{Your turn}
\begin{itemize}
\raggedright
  \item \emph{Say it:} \emph{samandar}
  \item \emph{Say it:} \emph{samandar}, once more
  \item \emph{Say it:} say \emph{samandar}
  \item \emph{Recall:} say the Marwadi for the sun, then the Marwadi for the moon, then say \emph{samandar} again
\end{itemize}

\begin{tcolorbox}[breakable,colback=teal!4,colframe=teal!35!black,title={Before you move on}]
What does \mw{समंदर} mean? (\textquotedblleft{}The sea\textquotedblright{}.) Say it once more.
\end{tcolorbox}
\section[\texorpdfstring{\mw{आग} (āg)}{आग (āg)}]{\mw{आग} (āg) — fire}

\label{lesson:MW-C69-ag}



\begin{quote}
Before the new one: say the Marwadi for the moon, then the Marwadi for the sea.
\end{quote}

\subsection*{You'll want to know: \mw{आग}}
\textbf{\mw{आग}} — \emph{āg} — \textquotedblleft{}fire\textquotedblright{}.

\begin{cousinweb}
\textbf{\mw{आग}} (\emph{āg}) is a tadbhava, a word inherited from Sanskrit and reshaped by sound change. It comes from \emph{agni}, \textquotedblleft{}fire\textquotedblright{}. CDIAL lists Old Marwari \emph{āgi} under it (CDIAL 55). Latin \emph{ignis}, \textquotedblleft{}fire\textquotedblright{}, is the same ancient word.
\end{cousinweb}

\begin{grammarlens}[title={What you've built}]
One more word for this chapter.
\end{grammarlens}

\subsection*{Your turn}
\begin{itemize}
\raggedright
  \item \emph{Say it:} \emph{āg}
  \item \emph{Say it:} \emph{āg}, once more
  \item \emph{Say it:} say \emph{āg}
  \item \emph{Recall:} say the Marwadi for the moon, then the Marwadi for the sea, then say \emph{āg} again
\end{itemize}

\begin{tcolorbox}[breakable,colback=teal!4,colframe=teal!35!black,title={Before you move on}]
What does \mw{आग} mean? (\textquotedblleft{}Fire\textquotedblright{}.) Say it once more.
\end{tcolorbox}
